\scanpage{158}
\mspage{32}
% archive sheet number: 157

Ces propriétés s'explicitent ainsi :

\uline{M1} Propriétés de relèvement-prolongement
\[
C\unsure{\longleftarrow|\longleftrightarrow}TF\quad\text{et}\quad TC\longleftrightarrow F
\]

\uline{M2} Propriété de factorisation de \deleted{toute} flèche
$f$ de $M$ en $f=pi$, avec, au choix, $i\in C,p\in TF$
ou $i\in TC,\ p\in F$.

\[
M'\qquad TC\subset C\cap W,\quad TF\subset F\cap W
\]
\deleted{\uline{Ex} $W=TF\circ TC$ (\unsure{axiome de saturation})}

C'est donc une partie des axiomes d'une
str. de Quillen faible \deleted{\unsure{aussi}} pour $W$
($\overline{\underline{\mathrm{XIII}}}$ p. 27 déf. 1) -- il manque les axiomes
de stabilité M3, M4, et les axiomes de
``propriété'' (\unsure{propriety}) M$''$, savoir $F\subset\mathrm{Fib}_W$, $C\subset\mathrm{Cof}_W$.

\uline{Proposition 4.3.} Soient $M$ une catégorie, $(C,TF)$ et
$(TC,F)$ deux \inserted{pré-}couples de Quillen dans $M$, satisfaisant
\[
TC\subset C,\quad TF\subset F
\tag*{(4.3.1.)}
\]
et soit $W$ une partie de $\mathrm{Fl}(M)$ dont on ne
suppose pour l'instant que 4.6. b) *, et
% margin note written sideways
\margin{* On utilise seulement que si $x\overset f{\unsure{-|\to}}y\overset g{\unsure{-|\to}}z$, et si $gf\in W$, alors $f\in W\Longleftrightarrow g\in W$. On n'a pas besoin que $W$ contienne les iso, ni qu'il soit stable par composition.}
\[
W\supset TC,TF,\quad\text{donc }\deleted{\unsure{F}}\text{on a}
\tag*{(4.3.2)}
\]
\[
TC\subset C\cap W,\quad TF\subset F\cap W.
\tag*{(4.3.3)}
\]
\deleted{C'est l'axiome de propriété de saturation pour $W$).}
Considérons les conditions de \deleted{factorisation} \inserted{prolongement-relèvement} suivantes
a), a$'$), et les conditions de factorisation b) b$'$) \deleted{\unsure{b}}c), d)
(en se rappelant que $\widetilde{TC}=F^*$, $\overline{TF}=C_*$, en
vertu de prop. 4.2 (cor. 1) appliqué aux \inserted{pré-}couples de
Quillen $(TC,F)$, $(C,TF)$) :

\[
\begin{aligned}
\text{a)}\quad &\deleted{\unsure{F}}\ C\cap W\subset\widetilde{TC}
&&(\text{donc }TC\subset C\cap W\subset\widetilde{TC})\\
\text{a$'$)}\quad &F\cap W\subset\overline{TF}
&&(\text{donc }TF\subset F\cap W\subset\overline{TF})\\
\text{b)}\quad &W\subset\deleted{\overline{\vphantom{TF}}}TF\circ\widetilde{TC}\\
\text{b$'$)}\quad &W\subset\overline{TF}\circ\deleted{\widetilde{\vphantom{TC}}}TC\\
\text{c)}\quad &W\subset\overline{TF}\circ\widetilde{TC}\\
\text{d)}\quad &W\subset TF\circ TC&&(\text{donc }\inserted{\text{par (4.3.2.)}}\ W=TF\circ TC\text{ si }W\text{ est stable}\\
&&&\text{par composition}
\end{aligned}
\]
% slanted
\margin{\uline{NB} \deleted{\unsure{aux}} a) signifie \unsure{juste|aussi} que $(C\cap W,F)$ est encore un précouple de Quillen, a\unsure{$'$)|)} signifie que $(C,F\cap W)$ en est un -- et alors cette paire de précouples de Quillen satisfait cette fois à d) i.e. $W\unsure{\subset|=}TF'\circ TC'$}
